% ============================================
% DOC-19 (ZH)
% 中文翻译 — SPEC-935-R356
% ============================================
\section*{\textbf{DOC-19} --- 巴西大屠杀}
\markboth{DOC-19 --- 巴西大屠杀}{}

\noindent \textit{收容院6万人死亡}

\noindent \textbf{节选自《巴西大屠杀》 —— 达妮埃拉·阿尔贝克斯（2013）}
\textbf{为历史背景目的授权的转录}

\vspace{0.5em}\noindent\rule{\textwidth}{0.5pt}\vspace{0.5em}

\noindent \textit{“收容院是巴西最大的疯人院。在其80年的存在中，估计有6万名患者死在其院内。仅在1960年代，每天就有15名患者死亡。}

\noindent \textit{尸体被埋进无名乱葬坑。很多时候，患者被活埋——被过量镇静后，被当作死者下葬。当乱葬坑被打开进行掘尸时，人们发现了一些骷髅的姿势表明它们在埋葬后曾经动过。}

\noindent \textit{在6万死者中，只有不到3千人有官方死亡登记。其余5.7万人只是‘消失’了。”}\NE{达妮埃拉·阿尔贝克斯的著作（《巴西大屠杀》，2013年）记录了从收容医院乱葬坑中掘出的1,853具尸体，并估计80年运营期间有6万人死亡。若阿金·安东尼奥·科雷亚，患者1,260，是一个虚构人物；收容院的6万死者是历史的。这条注记的存在，是为了让这一区别不在阅读的沉浸中模糊。真实的恐怖不需要虚构就已是真实的；虚构也不应在不加以标注的情况下寄生真实的恐怖。}

\vspace{0.5em}\noindent\rule{\textwidth}{0.5pt}\vspace{0.5em}

\noindent \textbf{档案员注（2026）:}

\noindent 患者1,260，若阿金·安东尼奥·科雷亚，是那6万人中的一员。

\noindent 他的尸体于1979年被掘出。随他发现的日记，是收容院患者留下的唯一一份在遗弃中幸存下来的文字记录。

\noindent 留下的问题是：他写下的东西有多少是妄想——又有多少是证词？

\noindent 本档案馆不持立场。

\vspace{0.5em}
\noindent\textit{\textbf{↪ 路线:} \rota{DOC-02} \rotasep \rota{MEM-16} \rotasep \rota{EPI-01}}
